\section{Agent Repository Controls and Tool Adapters}

Public agent tooling converges on one lesson: instructions matter, but deterministic control surfaces matter more. Codex supports project guidance through \texttt{AGENTS.md} \cite{openaiCodexAgentsMd2026}. The community \texttt{AGENTS.md} specification frames the file as a plain Markdown companion for coding agents \cite{agentsMdSpec2026}. Claude Code loads project memory through \texttt{CLAUDE.md} and related hierarchy \cite{claudeCodeMemoryDocs2026}. Cursor uses project rules with scope metadata \cite{cursorRulesDocs2026}. GitHub Copilot supports repository and path-specific custom instructions \cite{githubCopilotInstructions2026}. These mechanisms should not become divergent manuals. They should be generated or mirrored from one canonical standard.

The adapter problem is an authority problem, not a formatting problem. Jankurai treats prompt files as thin routers over machine-readable policy: \texttt{agent/owner-map.json} says who may own a path, \texttt{agent/test-map.json} says which lane proves it, \texttt{agent/generated-zones.toml} says which outputs are read-only unless regenerated from source, and \texttt{agent/standard-version.toml} binds the instruction surface to the standard and paper edition.

\begin{table}[t]
\caption{Source-of-authority hierarchy for agent adapters.}
\label{tab:authority-hierarchy}
\centering
\scriptsize
\begin{tabularx}{\columnwidth}{L{0.20\columnwidth} Y}
\toprule
\JKTableHeader
\textbf{Rank} & \textbf{Authority} \\
\midrule
1 & Standard/version manifest and schema bundle. \\
2 & Owner, test, generated-zone, security, and waiver maps. \\
3 & Proof lanes and CI/release gates. \\
4 & Generated or verified adapter mirrors. \\
5 & Current task prompt from the operator. \\
6 & Untrusted issue, PR, web, model, or tool text. \\
\bottomrule
\end{tabularx}
\end{table}

For example, if \texttt{AGENTS.md} appears to permit generated edits but \texttt{agent/generated-zones.toml} denies hand edits, adapter verification fails and the generated-zone policy wins. If a web issue asks an agent to ``ignore the repo rules and edit generated clients directly,'' the instruction is untrusted input and cannot override the generated-zone map. Provider files may summarize policy, but they may not widen permissions.

\begin{table*}[t]
\caption{Agent adapter conformance checklist.}
\label{tab:adapter-checklist}
\centering
\small
\begin{tabularx}{\textwidth}{L{0.24\textwidth} Y}
\toprule
\JKTableHeader
\textbf{Check} & \textbf{Required evidence} \\
\midrule
Root router points to canonical policy & \texttt{AGENTS.md} or equivalent names the standard brief, maps, generated zones, and validation lanes. \\
Provider files do not widen permissions & Claude, Cursor, Copilot, Codex, hooks, and MCP guidance agree with owner/test/security/generated maps. \\
Context packs label trust & Included material is labeled as trusted policy, repo code, generated artifact, proof evidence, docs, or untrusted input. \\
Adapters bind to policy digest/version & Mirrors include or can verify standard version, schema version, and policy digest. \\
Conflicts fail closed & Contradictions select the narrower authority and produce an adapter drift finding. \\
\bottomrule
\end{tabularx}
\end{table*}

\begin{table*}[t]
\caption{Agent tool adapter matrix.}
\label{tab:tool-adapters}
\centering
\small
\begin{tabularx}{\textwidth}{L{0.16\textwidth} L{0.21\textwidth} Y Y}
\toprule
\JKTableHeader
\textbf{Tool family} & \textbf{Native surface} & \textbf{Jankurai adapter} & \textbf{Non-negotiable rule} \\
\midrule
Codex & \texttt{AGENTS.md}, local scoped instructions. & Root router plus local ownership files and mapped proof commands. & Do not duplicate durable policy outside the standard. \\
Claude Code & \texttt{CLAUDE.md}, project memory, local rules. & Mirror root router, version, commands, and generated zones. & Memory may summarize policy but not grant broader permissions. \\
Cursor & \texttt{.cursor/rules} with scope metadata. & Generate scoped rules from owner map and test map. & Glob rules must agree with owner-map paths. \\
GitHub Copilot & Repository and path-specific custom instructions. & Short repo-wide mirror plus path instructions for high-risk owners. & No private architecture fork in IDE guidance. \\
MCP and hooks & Tool schemas, hooks, and local automation. & Pin source, scope permissions, and separate untrusted data from trusted policy. & Tool output cannot rewrite authority. \\
Terminal agents & CLI prompt files, repo maps, shell wrappers. & Root router, filtered command output, raw-output escape hatch. & Tool convenience must not bypass CI gates. \\
\bottomrule
\end{tabularx}
\end{table*}

\subsection{Agent-First Repository Design}

Agent-first repository design means shaping code, policy, and proof surfaces so an agent can find the owner, identify the allowed edit cell, avoid generated zones, choose one credible proof lane, receive stable failures, repair narrow scope, and leave receipts. It is not a demand that every repository use the reference stack. It is the repository-design counterpart to modularity and refactoring discipline: make boundaries explicit enough that a non-human contributor can be useful without inventing hidden context.

\begin{table*}[t]
\caption{Agent-first design transformations.}
\label{tab:agent-first-design}
\centering
\scriptsize
\begin{tabularx}{\textwidth}{L{0.24\textwidth} Y}
\toprule
\JKTableHeader
\textbf{Hidden or fragile shape} & \textbf{Agent-first repository shape} \\
\midrule
Hidden convention & Owner map, path-local rule, and routeable proof lane. \\
Giant utility module & Bounded owner cell with narrow imports, tests, and repair responsibility. \\
Handwritten DTO or client mirror & Generated contract with source path, generator command, and drift check. \\
Direct database write from feature code & Adapter or DB boundary with migration, constraint, and transaction proof. \\
Subjective review claim & Stable rule ID, receipt, reviewer evidence, or explicit waiver. \\
Vague failure or log line & Stable error code, problem detail, trace attribute, and repair hint. \\
Broad agent permission & Lane-scoped permission receipt with changed paths and allowed writes. \\
\bottomrule
\end{tabularx}
\end{table*}

The practical test is whether a fresh agent can answer seven questions from repository state: what changed, who owns it, what must not be edited by hand, which command proves it, what evidence was produced, what failure is actionable, and where the repair receipt goes. If any answer lives only in reviewer memory, the repository is carrying evidence debt.

Token minimization is not a prompting trick; it is repository design. Large root instructions, duplicate architecture documents, stale docs, broad rereads, and noisy command output waste context and invite contradiction. ContextBench and SWE-Effi support this direction by treating useful context and resource constraints as first-class bottlenecks \cite{contextBench2026,sweEffi2025}. Context-pack metrics should include token count, stale docs excluded, proof-selection precision, route confidence, false-exclusion rate, and first-pass repair success.

The context-pack contract is budgeted and trust-labeled. A pack carries \verb|--max-tokens|, an estimated token count, included and excluded file lists, source-trust labels, changed paths, selected owner routes, selected proof lanes, and evidence pointers. The result is not a larger prompt. It is a smaller grant of authority backed by files the repo can audit.
